%%%%%%%%%%%%%%%%%%%%%%%%%%%%%%%%%%%%%%%%%%%%%%%%%%%%%%%%%%%%%%%%%%%%%%%%%%%
\documentclass{beamer}

\hypersetup{pdfpagemode=FullScreen} %full screen mode
\setbeamertemplate{navigation symbols}{}
\usepackage[slovene]{babel}
%\usepackage[english]{babel} %designed for typesetting lebels(addr lebel,sticky lebels,etc)

\usepackage[utf8]{inputenc}
%\usepackage{lmodern}  % hm, če to importiraš, grške črke ne delajo??

\def\qed{$\hfill\Box$}   % konec dokaza
\def\qedm{\qquad\Box}   % konec dokaza v matematičnem načinu

%\usepackage[latin1]{inputenc} %Accept different input encodings
\usepackage{amsmath,amssymb, amsmath, mathtools} 
\usepackage{graphicx} % required for inserting images
\usepackage[T1]{fontenc} %for font encoding
\usepackage{times} % use times font instead of default
\usepackage{curves} %for drawing curves
\usepackage{verbatim} %paragraph making environment 
\usepackage{multimedia} %for multimedia like animation,movie etc...
\usepackage{mathptmx} % font style
\usepackage{graphicx} % Allows including images
\usepackage{booktabs} % Allows the use of \toprule, \midrule and \bottomrule in tables
\usepackage{hyperref}
\usepackage{xcolor}
\usepackage[symbol]{footmisc}
\usepackage{mathtools}
\usepackage{tikz} 
\usepackage{multirow}
\usepackage{qtree}
\usepackage{stmaryrd}

\usepackage{calligra}
\usepackage[T1]{fontenc}

\usepackage{algorithm,algorithmic}
\usepackage{ulem}

\makeatletter
\def\bign#1{\mathclose{\hbox{$\left#1\vbox to8.5\p@{}\right.\n@space$}}\mathopen{}}
\def\Bign#1{\mathclose{\hbox{$\left#1\vbox to11.5\p@{}\right.\n@space$}}\mathopen{}}
\makeatother

\usepackage{proof}
\usepackage{dsfont} 

\DeclareMathOperator{\Sym}{Sym}
\DeclareMathOperator{\Perm}{Perm}
\DeclareMathOperator{\supp}{supp}
\DeclareMathOperator{\Terms}{Terms}
\DeclareMathOperator{\sub}{sub}
\DeclareMathOperator{\vars}{vars}

\newcommand{\inference}[3][]{\infer[#1]{#3}{#2}}

\def\N{\mathbb{N}} % mnozica naravnih stevil
\def\Z{\mathbb{Z}} % mnozica celih stevil
\def\Q{\mathbb{Q}} % mnozica racionalnih stevil
\def\R{\mathbb{R}} % mnozica realnih stevil
\def\C{\mathbb{C}} % mnozica kompleksnih stevil

% \DeclareMathOperator{\ord}{ord}

\renewcommand{\algorithmicrequire}{\textbf{Input:}}
\renewcommand{\algorithmicensure}{\textbf{Output:}}
\usepackage{lipsum}
\setbeamertemplate{caption}[numbered] % For numbering figures

\def\qed{$\hfill\Box$}   % konec dokaza
\newtheorem{izrek}{Izrek}
\newtheorem{trditev}{Trditev}
\newtheorem{posledica}{Posledica}
\newtheorem{lema}{Lema}
\newtheorem{definicija}{Definicija}
\newtheorem{pripomba}{Pripomba}
\newtheorem{primer}{Primer}
\newtheorem{zgled}{Zgled}
% \newtheorem{zgledi}{Zgledi uporabe}
% \newtheorem{zglediaf}{Zgledi aritmetičnih funkcij}
% \newtheorem{oznaka}{Oznaka}
\newtheorem{domneva}{Domneva}
\newtheorem{dokaz}{Dokaz}
\newtheorem{opomba}{Opomba}
\newtheorem{oznaka}{Oznaka}
\newtheorem{izrek_crosser}{Izrek (Church–Rosserjev izrek)}

\mode<presentation> {

% The Beamer class comes with a number of default slide themes
% which change the colors and layouts of slides. Below this is a list
% of all the themes, uncomment each in turn to see what they look like.

%\usetheme{default}
%\usetheme{AnnArbor}
%\usetheme{Antibes}
%\usetheme{Bergen}
%\usetheme{Berkeley}
%\usetheme{Berlin}
%\usetheme{Boadilla}
%\usetheme{CambridgeUS}
%\usetheme{Copenhagen}
%\usetheme{Darmstadt}
%\usetheme{Dresden}
%\usetheme{Frankfurt}
%\usetheme{Goettingen}
%\usetheme{Hannover}
%\usetheme{Ilmenau}
%\usetheme{JuanLesPins}
%\usetheme{Luebeck}
\usetheme{Madrid}
%\usetheme{Malmoe}
%\usetheme{Marburg}
%\usetheme{Montpellier}
%\usetheme{PaloAlto}
%\usetheme{Pittsburgh}
%\usetheme{Rochester}
%\usetheme{Singapore}
%\usetheme{Szeged}
%\usetheme{Warsaw}

% As well as themes, the Beamer class has a number of color themes
% for any slide theme. Uncomment each of these in turn to see how it
% changes the colors of your current slide theme.

%\usecolortheme{albatross}
%\usecolortheme{beaver}%this one also
%\usecolortheme{beetle}
%\usecolortheme{crane} % also use this one
%\usecolortheme{dolphin}
%\usecolortheme{dove}
%\usecolortheme{fly}
%\usecolortheme{lily}
%\usecolortheme{orchid}
%\usecolortheme{rose}
%\usecolortheme{seagull}
%\usecolortheme{seahorse}
\usecolortheme{whale} % Best one
%\usecolortheme{wolverine} %use can use this also

%\setbeamertemplate{footline} % To remove the footer line in all slides uncomment this line
%\setbeamertemplate{footline}[page number] % To replace the footer line in all slides with a simple slide count uncomment this line

%\setbeamertemplate{navigation symbols}{} % To remove the navigation symbols from the bottom of all slides uncomment this line
}

\usepackage{graphicx} % Allows including images
\usepackage{booktabs} % Allows the use of \toprule, \midrule and \bottomrule in tables
\setbeamercovered{transparent}
\setbeamertemplate{bibliography item}[text]
\setbeamertemplate{theorems}[numbered]
\setbeamerfont{title}{size=\Large}%\miniscule,command,tiny, scriptsize,footnotesize,small,normalsize,large,Large,LARGE,huge,Huge,HUGE
\setbeamerfont{date}{size=\tiny}%{\fontsize{40}{48} \selectfont Text}

\setbeamertemplate{itemize items}[triangle]
\setbeamertemplate{itemize subitem}[circle]
\setbeamertemplate{itemize subsubitem}[ball] 


%------------------customized frame----------------------------
\newcounter{cont}

\makeatletter
%allowframebreaks numbering in the title
\setbeamertemplate{frametitle continuation}{%
   % \setcounter{cont}{\beamer@endpageofframe}%
    %\addtocounter{cont}{1}%
   % \addtocounter{cont}{-\beamer@startpageofframe}%
   % (\insertcontinuationcount/\arabic{cont})%
}
\makeatother
%-----------------------customized-----------------------------




%----------------------------------------------------------------------------------------
%	TITLE PAGE
%----------------------------------------------------------------------------------------

\title{Lambda izrazi v nominalnih množicah} % The short title appears at the bottom of every slide, the full title is only on the title page

\title[Lambda izrazi v nominalnih množicah]{Lambda izrazi v nominalnih množicah}  % sj je isto
\author[Tadej Glinšek]{Tadej Glinšek\\{\small Mentor: prof.~dr.~Alexander Keith Simpson}}
\institute[]{Fakulteta za matematiko in fiziko \\
Oddelek za matematiko}



 % Your institution as it will appear on the bottom of every slide, may be shorthand to save space
{
%\institute[FMF] % Your institution as it will appear on the bottom of every slide, may be shorthand to save space


\medskip
%\begin{center}
%\includegraphics[width=0.1\textwidth]{uoh.png}
%\end{center}
\vspace{2cm}
\medskip
% Fakulteta za matematiko in fiziko% Your institution for the title page



%\medskip
%\textit{john@smith.com} % Your email address
}
%This will place the image at position "30 right/left and 120 up/down" relative to the top left corner of the current page.

 \date{4.\ 5.\ 2026} % Date, can be changed to a custom date

\begin{document}

\begin{frame}
\titlepage % Print the title page as the first slide
\end{frame}
% \maketitle?

%\begin{frame}
  %\frametitle{Content of the presentation}
 % \tableofcontents[pausesections,shaded]
 %\tableofcontents
%\end{frame}

% TABLE OF CONTENTS AT BEGIN OF EACH SECTION
%----------------------------------------------------------------------------------------
%	PRESENTATION SLIDES
%----------------------------------------------------------------------------------------


%-------------------------------------------------------------------------------
%-------------------------------INTRODUCTION--------------------------------
%-----------------------------------------------------------


\begin{frame}[fragile]
\frametitle{Netipiziran lambda račun}

Netipiziran lambda račun:
\begin{itemize}
\item Alonzo Church, 1930-ta
\item \(\Lambda := \text{množica lambda izrazov}\)
\item Množico \(\lambda\)-izrazov lahko opišemo z gramatiko:\[
    t \in \Lambda ::= a \ \vert \ \lambda a.\ t\ \vert \ t\ t,
\]kjer je \(a \in \vars\) spremenljivka.
\item spremenljivka, abstrakcija, aplikacija
\end{itemize}
\pause
\begin{definicija}
    \(\lambda\)-izrazi v netipiziranem lambda računu so definirani z naslednjimi pravili:
    \[
        \inference[(x \in \vars)]{\quad \quad \quad \quad \quad}{x}
    \]
    \[
        \inference{M \quad \quad \quad N}{M N}\quad \quad \quad \quad
    \]
    \[
        \inference[(x \in \vars)]{\quad \quad M \quad \quad}{\lambda x.\ M}    
    \]
\end{definicija}

\end{frame}


\begin{frame}[fragile]
\frametitle{Predpogoji}

\begin{definicija}
    Naj bo \(G\) grupa. Množica \(X\) je \textbf{\(G\)-množica}, če grupa \(G\) deluje na množici \(X\), torej, za vsak \(g_1, g_2 \in G\) in \(x \in X\) velja
    \[
    (g_1 \cdot g_2) \cdot x = g_1 \cdot (g_2 \cdot x)
    \]
    in za enoto \(e\) velja
    \[
        e \cdot x = x.
    \]
\end{definicija}

\pause
\begin{definicija}
Naj bo \(\mathds{A}\) neskončna množica in \(\Sym(\mathds{A})\) grupa vseh permutacij na \(\mathds{A}\). Permutacija \(\pi \in \Sym(\mathds{A})\) je \textbf{končna}, če je \(\{a \in \mathds{A}\ |\ \pi a \neq a\}\) končna množica.
\end{definicija}

\begin{definicija}
Množico končnih permutacij iz \(\Sym(\mathds{A})\) označimo z oznako \(\Perm(\mathds{A})\).    
\end{definicija}


\end{frame}


%%%%%%%%%%%%%%%%%%%%%%%%%%%%%%%%%%%%%%%%%%%%%%%%%%%%%%%%%


\begin{frame}[fragile]
\frametitle{Nosilec}
\begin{definicija}

\textbf{Nosilec} (angl. \textbf{support}) elementa \(x \in X\) je množica \(A \subseteq \mathds{A}\), da za vsak \(\pi \in \Perm \mathds{A}\) velja
\[
    (\forall a \in A.\ \pi\ a = a) \Rightarrow \pi \cdot x = x.
\]
\end{definicija}
\begin{opomba}
    Nosilcev elementa \(x\) je lahko več.
\end{opomba}
\pause

\begin{trditev}
    Množica \(A \subseteq \mathds{A}\) je nosilec za \(x \in X\) natanko tedaj, ko 
    \[
        \forall a,\ a' \in \mathds{A}\setminus A.\ (a\ a') \cdot x = x
    \]
\end{trditev}


\end{frame}






\begin{frame}[fragile]
\frametitle{Končni nosilci}


\begin{definicija}
    \textbf{Nominalna množica} je \(\Perm\mathds{A}\)-množica, katere vsak element ima končen nosilec.
\end{definicija}
\pause
\begin{lema}
    Naj bosta \(A\) in \(A'\) končna nosilca elementa \(x\). Potem je \(A\, \cap\, A'\) spet nosilec.
\end{lema}

\begin{opomba}
    Sledi, da če obstaja končen nosilec elementa \(x\), potem obstaja najmanjši nosilec. Označimo ga \(\supp_X(x)\).
\end{opomba}

\end{frame}


\begin{frame}[fragile]
\frametitle{Produkti}

\begin{definicija}
    Naj bodo \(X_i\), kjer \(i = 1,\ \ldots,\ n\), nominalne množice. Potem je \(X_1 \times \cdots \times X_n\) spet nominalna množica:
    \[
        \pi \cdot (x_1,\ \ldots,\ x_n) = (\pi \cdot x_1,\ \ldots,\ \pi \cdot x_n)
    \]
\end{definicija}
\pause
\begin{trditev}
    Za nominalne množice \(X_i\), kjer \(i = 1,\ \ldots,\ n\), velja:
    \[
        \supp_{X_1 \times \cdots \times X_n}(x_1,\ \ldots,\ x_n) = \supp_{X_1} x_1\, \cup\, \cdots\, \cup\, \supp_{X_n} x_n
    \]
\end{trditev}


\end{frame}
%%%%%%%%%%%%%%%%%%%%%%%%%%%%%%%%%%%%%%%%%%%%%%%%%%%%%%%%%

\begin{frame}[fragile]
\frametitle{Relacije}
\begin{definicija}
    Naj bo \(X\) nominalna množica in \(x, y \in X\) elementa.
    Elementa \(x\) in \(y\) sta \textbf{orbitno ekvivalentna} (\(x \sim y\)), če velja, da obstaja \(\pi \in \Perm \mathds{A}\), da
    \[
        \pi \cdot x = y
    \]
    
\end{definicija}
\pause
\begin{definicija}
    Naj bosta \(X, Y\) nominalni množici, \(x \in X\), \(y \in Y\). Pravimo, da sta \(x\) in \(y\) \textbf{separirana} (\(x \ \#\ y\)), če velja:
    \[
        x \ \#\ y\ :\Longleftrightarrow\ \supp x\, \cap\, \supp y = \emptyset
    \]
\end{definicija}
    
\end{frame}


\begin{frame}[fragile]
\frametitle{Relacije}

\begin{definicija}
    Naj bodo \(X, Y, Z\) nominalne množice, \(x \in X\), \(y \in Y\), \(z \in Z\). Pravimo, da sta \(x\) in \(y\) \textbf{pogojno separirana} glede na \(z\), če velja
    \[
        x \ \#\  y \ \vert \ z\ :\Longleftrightarrow\ \supp(x)\, \cap\, \supp(y) \subseteq \supp(z)
    \]
\end{definicija}

\end{frame}

\begin{frame}[fragile]
\frametitle{Lastnosti relacij}

\begin{trditev}
    Naj bosta \(X\) in \(Y \subseteq X\) nominalni množici. Naj \(x \in X\) in \(y \in Y\). Potem \(x \sim y \Rightarrow x \in Y\).
\end{trditev}
\pause
\begin{trditev}
    Za vsaka \(x,\ y \in X\) obstaja \(z \in X\), da velja \(z \ \#\  x\) in \(z \sim y\).
\end{trditev}
\pause
\begin{trditev}
    Za vsako trojico \(x,\ y,\ w \in X\) obstaja \(z \in X\), da velja \(z \ \#\  x\ \ \vert\ \ w\) in \(z \sim y\ \ \vert\ \ w\).
\end{trditev}

\end{frame}

\begin{frame}[fragile]
\frametitle{Lastnosti relacij}

\begin{trditev}
    Naj bodo \(X,\ Y,\ Z,\ W\) nominalne množice, \(x,\ y,\ z,\ w\) pa njihovi elementi. Potem velja:
    \[
        x \ \#\  (y,\ z) \ \vert\ w\ \Longleftrightarrow\ x \ \#\  y \ \vert\ (z,\ w)\, \wedge\, x \ \#\ z \ \vert\ w
    \]
\end{trditev}
\pause
\begin{zgled}
    Naj bodo \(A,\ B,\ C\) slučajne spremenljivke. Če velja
    \[
        A\ \perp\ B\ \vert\ C\ :\Longleftrightarrow\ P((A,\ B) \ \vert\ C) = P(A\ \vert\ C)P(B\ \vert\ C)
    \] pravimo, da sta \(A\) in \(B\) pogojno neodvisni glede na \(C\).
    
    Naj bodo \(X,\ Y,\ Z,\ W\) slučajne spremenljivke. Potem velja
    \[
        X \, \perp\,  (Y,\ Z) \ \vert\ W\ \Longleftrightarrow\ X \, \perp\,  Y \ \vert\ (Z,\ W)\ \wedge\ X \, \perp\, Z \ \vert\ W
    \]
\end{zgled}
\end{frame}



\begin{frame}[fragile]
\frametitle{Lambda izrazi nad kontekstom}

% \(\lambda\)-izraze definiramo skupaj s kontekstom:

\begin{definicija}
    \(\lambda\)-izraz \(M\) nad kontekstom \(\Gamma\) je definiran z naslednjimi pravili:
    \[
        \inference[(x \in \Gamma)]{}{\Gamma \vdash x}
    \]
    \[
        \inference{\Gamma \vdash M \quad \Gamma \vdash N}{\Gamma \vdash M N}
    \]
    \[
        \inference[(x \notin \Gamma)]{\Gamma,\ x \vdash M}{\Gamma \vdash \lambda x.\ M}    
    \]
Kontekst \(\Gamma\) definiramo kot končno zaporedje elementov množice \(\mathds{A}\).
\end{definicija}
\pause
\begin{definicija}
\[
    \Terms_{\Gamma} := \{\Gamma \vdash M\ |\ \Gamma \vdash M \ \text{je \(\lambda\)-izraz}\}
\]
\[
    \Terms := \{\Gamma \vdash M\ |\ \Gamma\ \text{je kontekst},\ \Gamma \vdash M \ \text{je \(\lambda\)-izraz}\}
\]
\end{definicija}
\end{frame}

% za kontekst nekako velja, da vsebuje vse proste spremenljivke


\begin{frame}[fragile]
\frametitle{Delovanje permutacij na množico \(\Terms\)}

\begin{definicija}
    Naj bo \(\Gamma = a_1,\,\ a_2,\ \ldots,\ a_n\) končno zaporedje elementov iz \(\mathds{A}\). Permutacija \(\pi \in \Perm\mathds{A}\) slika zaporedje \(\Gamma\) v \(\pi(a_1), \ \pi(a_2), \ \ldots, \ \pi(a_n)\).
\end{definicija}
\pause
\begin{definicija}
    Delovanje grupe \(\Perm{\mathds{A}}\) na množico izrazov je definirano sledeče:
    \[
        \pi(\Gamma \vdash x) := \pi(\Gamma) \vdash \pi(x)
    \]
    \[
        \pi(\Gamma \vdash MM') := \pi(\Gamma ) \vdash \pi(M)\pi(M')
    \]
    \[
        \pi(\Gamma \vdash \lambda x.\ M) := \pi(\Gamma) \vdash \lambda \pi(x). \pi(M)
    \]
\end{definicija}



\end{frame}

\begin{frame}[fragile]
\frametitle{Substitucija}


\begin{definicija}
    Naj bosta \(\Gamma,\ x \vdash M\) in \(\Gamma \vdash N\) \(\lambda\)-izraza, tako da velja \[N \ \#\ M \ \vert\ \Gamma.\] Potem definiramo \(\Gamma \vdash M[N/x]\) kot
    
    \[z[N/x]\, :=\, \begin{cases*}
    N & če z = x,\\
    z & sicer,
    \end{cases*}
    \]
    \[
        (M'\ M'')[N/x]\, :=\, M'[N/x]\ M''[N/x],\ \ \ \ \ 
    \]
    \[
        (\lambda z.\ M')[N/x]\, :=\, \lambda z.\ (M'[N/x]).\quad \quad \quad
    \]
    
    
\end{definicija}
\pause
\begin{izrek}
    Če so zgornji pogoji izpolnjeni, je izraz \(\Gamma \vdash M[N/x]\) dobro definiran.
\end{izrek}
\end{frame}


%\iffalse
\begin{frame}[fragile]
\frametitle{test}

\begin{lema}
    Naj \(\Gamma \vdash M\) \(\lambda\)-izraz, \(z\) spremenljivka, in naj velja \(z \ \#\  M,\ \Gamma\). Potem velja \[
        \Gamma,\ z \vdash M.
    \]
\end{lema}

\end{frame}
%\fi


\begin{frame}[fragile]
\frametitle{\(\alpha\)-ekvivalenca}


\begin{definicija}
    Relacija \(\equiv_\alpha\) je definirana s sledečimi pravili:
    \[
        \inference[]{}{\Gamma,\ x\ \vdash x \ \equiv_\alpha\ x}
    \]
    \[
        \inference{\Gamma \vdash M\ \equiv_\alpha\ M' \quad \Gamma \vdash N\ \equiv_\alpha\ N'}{\Gamma \vdash MN\ \equiv_\alpha\ M'N'}
    \]
    \[
        \inference{\Gamma,\ x \vdash M\ \equiv_\alpha\ N}{\Gamma \vdash \lambda x.\ M\ \equiv_\alpha\ \lambda x.\ N}
    \]
    \[
        \inference{\Gamma \vdash M \quad \Gamma \vdash N \quad (\Gamma \vdash M) \sim (\Gamma \vdash N) \ \vert\ \Gamma}{\Gamma \vdash M\ \equiv_\alpha\ N}
    \]
    Če za dva izraza velja \(\Gamma \vdash M\ \equiv_\alpha\ N\), potem rečemo, da sta izraza \textbf{\(\alpha\)-ekvivalentna}.
\end{definicija}

\end{frame}


\begin{frame}[fragile]
\frametitle{\(\alpha\)-ekvivalenca}
\begin{trditev}
    Naj bo \(\Gamma \vdash M\) izraz nad kontekstom \(\Gamma\) in naj bo \(\Delta \vdash N\) izraz nad kontekstom \(\Delta\), tako da velja \((\Gamma \vdash M) \sim (\Delta \vdash N) \ \vert\ \Gamma\). Potem je \(\Gamma \vdash N\) dobro definiran izraz.
\end{trditev}
\end{frame}




\begin{frame}[fragile]
\frametitle{\(\beta\)-redukcija}

\begin{definicija}
    Relacija \(\beta\)-redukcija (označimo z \(\rightarrow_\beta\)) je definirana s sledečimi pravili:
    \[
        \inference{\Gamma \vdash M \rightarrow_\beta M' \quad \Gamma \vdash N}{\Gamma \vdash MN \rightarrow_\beta M'N}
    \]
    \[
        \inference{\Gamma \vdash N \rightarrow_\beta N' \quad \Gamma \vdash M}{\Gamma \vdash MN \rightarrow_\beta MN'}
    \]
    \[
        \inference{\Gamma,\ x \vdash M \rightarrow_\beta M'}{\Gamma \vdash \lambda x.\ M \rightarrow_\beta \lambda x.\ M'}
    \]
    \[
        \inference{\Gamma,\ x \vdash M \quad \Gamma \vdash N \quad N' \sim N \ \vert\ \Gamma \quad M \ \#\ N' \ \vert\ \Gamma}{\Gamma \vdash (\lambda x.\ M)N \rightarrow_\beta M[N'/x]}
    \]
\end{definicija}

\end{frame}


\begin{frame}[fragile]
\frametitle{\(\beta\)-redukcija}

\begin{lema}
    Naj \(\Gamma \vdash M \rightarrow_\beta M'\) in naj
    \(\Gamma \vdash M \equiv_\alpha N\).
    Potem velja, da obstaja \(N'\), da \(\Gamma \vdash M' \equiv_\alpha N' \) in \(\Gamma \vdash N \rightarrow_\beta N'\).
\end{lema}
\pause
\begin{lema}
    Naj bo \(\Gamma\) kontekst, \(x \in \mathds{A}\) pa se ne pojavi v \(\Gamma\). Naj velja \(\Gamma,\ x \vdash M,\ M'\) in \(\Gamma \vdash N,\ N'\), tako da \(\Gamma \vdash M \equiv_\alpha M'\) in \(\Gamma \vdash N \equiv_\alpha N'\). Naj velja še \(M \ \#\ N \ \vert\ \Gamma\) in \(M' \ \#\ N' \ \vert\ \Gamma\). Potem sta \(\Gamma \vdash M[N/x],\ M'[N'/x]\) dobro definirana in velja \(\Gamma \vdash M[N/x] \equiv_\alpha M'[N'/x]\).
\end{lema}


\end{frame}


\begin{frame}[fragile]
\frametitle{Church-Rosserjev izrek}
\begin{izrek_crosser}
    Naj bo \(\Gamma\) kontekst in \(\Gamma \vdash M,\ N_1,\ N_2\) \(\lambda\)-izrazi nad tem kontekstom. Naj \(\Gamma \vdash M \rightarrow_\beta N_1\) in \(\Gamma \vdash M \rightarrow_\beta N_2\). Potem obstaja \(\Gamma \vdash X\), da velja
    \[
        \Gamma \vdash N_1 \rightarrow_\beta X
    \]
    in
    \[
        \Gamma \vdash N_2 \rightarrow_\beta X.
    \]
\end{izrek_crosser}

\end{frame}
%--------------------------------------------------------------------------------------------------------------------------------
%\section{Information processing  issues}
%-----------------------------------------------
%\begin{frame}
%\frametitle{Information processing issues}
%\begin{itemize} 
%\item In collaborative processing
%\begin{itemize}
%\item Target detection, localization, communication, storage, querying...
%\end{itemize}
%\color{lightgray}
%\item  In networking
%\begin{itemize}
%\color{lightgray}
%\item  aggregation, routing, Data naming,...
%\end{itemize}
%\end{itemize}
%\end{frame}

%-----------------------------------------------
\iffalse
\section{References}
\begin{frame}[allowframebreaks]
\frametitle{References}
 \nocite{*}
 \bibliographystyle{unsrt}
\bibliography{References}
\end{frame}
\fi

%-----------------------------------------------------------------------------------------


%----------------------------------------------------------------------------------------

\end{document}
